%% GENERATED by python/paper/build.py from results/ -- do not edit by hand.
%% Rebuild:  .venv\Scripts\python.exe python\paper\build.py --only US_ESC_Ageing
%% Source:   results/esc/escExperiments.csv
\begin{table}[!htb]
\centering
\begin{threeparttable}
\caption{Endogenous design and ``acute ageing'' in US}
\label{table:US_ESC:ageing}
\renewcommand{\arraystretch}{1.25}
\begin{tabularx}{.9\textwidth}{p{2.6cm}YYYYY}
\toprule
\textbf{Scenario} & \textbf{CRRA} ($\rho$) & $\bm{\theta}$ \textbf{(2020)} & \textbf{Tax rate} & \textbf{Savings rate} & \textbf{Avg. workweek} \\
\midrule 
 & 0.5 & 0.74 & 14.42\% & 15.96\% & 39.42 \\
Baseline & 1.0 & 0.74 & 14.43\% & 15.38\% & 39.40 \\
 & 2.0 & 0.74 & 14.45\% & 14.84\% & 39.42 \\[.5em]\hline\\[-.75em]
 & 0.5 & 0.74 & 24.17\% & $-2.37$ p.p. & 38.61 \\
Exogenous $\theta$ & 1.0 & 0.74 & 23.25\% & $-2.09$ p.p. & 38.63 \\
 & 2.0 & 0.74 & 22.01\% & $-1.78$ p.p. & 38.71 \\[.5em]\hline\\[-.75em]
 & 0.5 & 0.77 & 24.04\% & $-2.41$ p.p. & 38.71 \\
Endogenous $\theta$ & 1.0 & 0.78 & 23.19\% & $-2.12$ p.p. & 38.69 \\
 & 2.0 & 0.85 & 22.06\% & $-1.83$ p.p. & 38.83 \\[.5em]\hline\\[-.75em]
\bottomrule
\end{tabularx}
\begin{tablenotes}
\footnotesize
\item Deadweight-cost specification: the proportional cost $f(\theta)$ with $\phi = 0.5$ and $p$ calibrated per $\rho$ (\cref{table:US_ESC:calibration}). Every counterfactual is a separate equilibrium path: the changed parameters hold throughout, the economy starts from its own steady state, and the political choice binds from the first period of the horizon, so the design in force in 2020 is itself an outcome rather than an inherited datum. All rows are read at 2020. $\theta$ (2020) is the design in force there; in the exogenous rows it is the US design, held fixed so that the counterfactual is about the changed characteristic alone. Each $\rho$ is separately calibrated and its workweek normalised against its own baseline. The savings rate is savings relative to GDP; the baseline rows report its level and every other row the change against the baseline at the same $\rho$, in percentage points. The ``acute ageing'' scenario sets $\nu_t = 1$ throughout, so the counterfactual economy is one whose demography has always been stationary rather than the US surprised by ageing in 2020. Common-$X$ calibration: one leisure parameter $X$ shared across income groups, with the hours unit pinned by targeting the observed average workweek, so relative hours are a prediction rather than a calibration target.
\end{tablenotes}
\end{threeparttable}
\end{table}
